\documentclass[11pt]{article}
\usepackage{mathblatt}
% gesetzt von werkzeuge/setzer.py aus dem Lernweg (L5g-A · L5g-B · L5g-C · L5g-D · L5g-E) und den Bankzeilen; Muster T6B.
% Präambel des Setzers (werkzeuge/setzer.py), wortgleich aus dem Hand-Satz T6B (bau/proben/2026-10-09/kathete-berechnen), dort aus K4W.
\makeatletter
\setlength{\mbpfbe}{0mm}% keine Punkte (Bauregel 6.4)
% eingerückt (Bauregel 4.7, 6.6): \gEin Einzug, \gSchrift eine Stufe kleiner
\newlength{\gEin}\setlength{\gEin}{0pt}
\let\gSchrift\relax
\renewcommand{\pfkreuzzeile}[1]{\par\noindent\hangindent3.2em\hangafter1\hspace*{1.6em}$\square$\ #1\par}
\newcommand{\gkreuz}[1]{$\square$\ #1\hspace{2em}}
% Bauregel 6.7: links, was man ansieht, rechts, was man tut; Spaltengrenze überall an derselben Stelle
\newsavebox{\gbox}\newlength{\gLinks}
\newcommand{\gzwei}[2]{\par\smallskip\noindent\sbox{\gbox}{#1}%
  \setlength{\gLinks}{\dimexpr0.5\textwidth-\gEin-\mbpfnr\relax}%
  \ifdim\wd\gbox>\dimexpr\gLinks-4mm\relax
    \usebox{\gbox}\par\smallskip\noindent\begin{minipage}{\linewidth}#2\end{minipage}%
  \else
    \begin{minipage}[t]{\gLinks}\vspace{0pt}\usebox{\gbox}\end{minipage}%
    \begin{minipage}[t]{\dimexpr\linewidth-\gLinks\relax}\vspace{0pt}\raggedright #2\end{minipage}%
  \fi\par}
\RenewDocumentEnvironment{pfaufg}{m m m}{%
  \begin{lrbox}{\mbaufgabebox}\begin{minipage}[t]{\linewidth}%
  \gSchrift\noindent\hspace*{\gEin}\mbpfrandmarke{#1}%
  \begin{minipage}[t]{\mbpfnr}\strut\textbf{#2}\end{minipage}%
  \begin{minipage}[t]{\dimexpr\linewidth-\gEin-\mbpfnr\relax}\raggedright\strut\ignorespaces}%
 {\end{minipage}%
  \end{minipage}\end{lrbox}%
  \setlength{\mbaufgabehoehe}{\dimexpr\ht\mbaufgabebox+\dp\mbaufgabebox\relax}%
  \par\addvspace{7pt}\Needspace*{\mbaufgabehoehe}%
  \noindent\usebox{\mbaufgabebox}\par\addvspace{7pt}}
\makeatother
% Rechenraum als Karo (Bauregel 6.9): n Rechenschritte = n mal 10 mm
\newcommand{\karo}[1]{\par\smallskip\noindent\begin{tikzpicture}
  \clip (0,0) rectangle ({\linewidth-1mm},{#1*10mm});
  \draw[black!16,line width=0.3pt,step=5mm] (0,0) grid ({\linewidth},{#1*10mm});
  \end{tikzpicture}\par}
\newcommand{\kk}[1]{$\square$\,#1\hspace{0.9em}}
\newcommand{\linie}[1]{\,{\color{black!45}\rule[-1pt]{#1}{0.4pt}}\,}
% Zeile am Ende jedes Durchgangs: Originale klein grau, darüber; dann Vorgänger/Nachfolger (Bauregel 3.4, 6.5)
\newcommand{\blattende}[2]{\par\vfill\noindent\begin{minipage}{\linewidth}{\scriptsize\color{mbgrau}Originale: #1\par}\smallskip
{\footnotesize #2\par}\end{minipage}}
\newcommand{\pkt}{\textperiodcentered{}}

\newcommand{\kennung}{FUR}
\newcommand{\grau}[1]{{\color{mbgrau}#1}}

\begin{document}
\pfheftstil
\fancyfoot[C]{\parbox[t]{\textwidth}{\mbfhbox\par\vspace{1.5mm}{\footnotesize\color{mbgrau}{\scriptsize \kennung}\hfill\thepage}}}
\pfheftkopf{Zahlen vergleichen in allen Darstellungen}{Klasse 8}
% L5g-A: Dezimalzahlen vergleichen
\begin{pfaufg}{}{1.}{}
Wo entscheidet es sich? Nenne nur die Stelle, an der sich die zwei Zahlen zuerst unterscheiden (Einer, Zehntel, Hundertstel, Tausendstel). Noch kein Zeichen setzen.\par\smallskip
\grau{a)\ $3{,}47$ und $3{,}42$:\quad Einer gleich, Zehntel gleich, Hundertstel verschieden $\Rightarrow$ Hundertstel}\par\medskip
\noindent\begin{minipage}[t]{0.48\linewidth}\raggedright
b)\ $0{,}81$ und $0{,}91$\linie{25mm}
\end{minipage}\hfill
\begin{minipage}[t]{0.48\linewidth}\raggedright
c)\ $12{,}305$ und $12{,}308$\linie{25mm}
\end{minipage}\par\smallskip
\noindent\begin{minipage}[t]{0.48\linewidth}\raggedright
d)\ $7{,}25$ und $6{,}25$\linie{25mm}
\end{minipage}\hfill
\begin{minipage}[t]{0.48\linewidth}\raggedright
e)\ $0{,}406$ und $0{,}46$\linie{25mm}
\end{minipage}\par\smallskip
\fusshilfe{\mbox{1:~b) Zehntel; c) Tausendstel; d) Einer; e) Hundertstel}}
\end{pfaufg}

% L5g-A: Dezimalzahlen vergleichen
\begin{pfaufg}{}{2.}{}
Hänge Nullen an, bis beide Zahlen gleich viele Stellen nach dem Komma haben. Setze dann $<$, $>$ oder $=$.\par\smallskip
\grau{a)\ $0{,}7$ \fbox{\phantom{$<$}} $0{,}65$:\quad $0{,}70$ und $0{,}65$\quad$\Rightarrow$\quad 7 Zehntel $>$ 6 Zehntel, also $0{,}7 > 0{,}65$}\par\medskip
\noindent\begin{minipage}[t]{0.48\linewidth}\raggedright
b)\ $2{,}4$ \fbox{\phantom{$<$}} $2{,}38$\linie{25mm}
\end{minipage}\hfill
\begin{minipage}[t]{0.48\linewidth}\raggedright
c)\ $0{,}09$ \fbox{\phantom{$<$}} $0{,}1$\linie{25mm}
\end{minipage}\par\smallskip
\noindent\begin{minipage}[t]{0.48\linewidth}\raggedright
d)\ $5{,}05$ \fbox{\phantom{$<$}} $5{,}5$\linie{25mm}
\end{minipage}\hfill
\begin{minipage}[t]{0.48\linewidth}\raggedright
e)\ $1{,}250$ \fbox{\phantom{$<$}} $1{,}25$\linie{25mm}
\end{minipage}\par\smallskip
\fusshilfe{\mbox{2:~b) $2{,}4 > 2{,}38$; c) $0{,}09 < 0{,}1$; d) $5{,}05 < 5{,}5$; e) $1{,}250 = 1{,}25$}}
\end{pfaufg}

% L5g-A: Dezimalzahlen vergleichen
\begin{pfaufg}{}{3.}{}
Ordne von klein nach groß: $0{,}6$; \ $0{,}06$; \ $0{,}066$; \ $0{,}61$.\karo{3}
\fusshilfe{\mbox{3:~$0{,}06 < 0{,}066 < 0{,}6 < 0{,}61$}}
\end{pfaufg}

% L5g-B: Negative Zahlen vergleichen
\begin{pfaufg}{}{4.}{}
Welche Zahl liegt auf der Zahlengeraden weiter links? Schreib sie auf. Die liegt weiter links, ist die kleinere.\par\smallskip
\grau{a)\ $-0{,}8$ und $-0{,}4$:\quad $-0{,}8$ liegt links von $-0{,}4$, also $-0{,}8 < -0{,}4$}\par\medskip
\noindent\begin{minipage}[t]{0.48\linewidth}\raggedright
b)\ $-2$ und $1$\linie{25mm}
\end{minipage}\hfill
\begin{minipage}[t]{0.48\linewidth}\raggedright
c)\ $-0{,}5$ und $0{,}3$\linie{25mm}
\end{minipage}\par\smallskip
\noindent\begin{minipage}[t]{0.48\linewidth}\raggedright
d)\ $-3$ und $-1$\linie{25mm}
\end{minipage}\hfill
\begin{minipage}[t]{0.48\linewidth}\raggedright
e)\ $0$ und $-0{,}1$\linie{25mm}
\end{minipage}\par\smallskip
\fusshilfe{\mbox{4:~b) $-2$; c) $-0{,}5$; d) $-3$; e) $-0{,}1$}}
\end{pfaufg}

% L5g-B: Negative Zahlen vergleichen
\begin{pfaufg}{}{5.}{}
Setze $<$ oder $>$. Bei zwei negativen Zahlen: Die weiter von 0 weg ist, ist die kleinere.\par\smallskip
\grau{a)\ $-0{,}6$ \fbox{\phantom{$<$}} $-0{,}65$:\quad $0{,}60 < 0{,}65$: \ $-0{,}65$ ist weiter von 0 weg\quad$\Rightarrow$\quad $-0{,}6 > -0{,}65$}\par\medskip
\noindent\begin{minipage}[t]{0.48\linewidth}\raggedright
b)\ $-1{,}2$ \fbox{\phantom{$<$}} $-1{,}02$\linie{25mm}
\end{minipage}\hfill
\begin{minipage}[t]{0.48\linewidth}\raggedright
c)\ $-0{,}3$ \fbox{\phantom{$<$}} $0{,}1$\linie{25mm}
\end{minipage}\par\smallskip
\noindent\begin{minipage}[t]{0.48\linewidth}\raggedright
d)\ $-2{,}5$ \fbox{\phantom{$<$}} $-2{,}45$\linie{25mm}
\end{minipage}\hfill
\begin{minipage}[t]{0.48\linewidth}\raggedright
e)\ $-0{,}09$ \fbox{\phantom{$<$}} $-0{,}1$\linie{25mm}
\end{minipage}\par\smallskip
\fusshilfe{\mbox{5:~b) $-1{,}2 < -1{,}02$; c) $-0{,}3 < 0{,}1$; d) $-2{,}5 < -2{,}45$; e) $-0{,}09 > -0{,}1$}}
\end{pfaufg}

% L5g-B: Negative Zahlen vergleichen
\begin{pfaufg}{}{6.}{}
Welche Zahl liegt genau in der Mitte? Rechne den Abstand aus, halbiere ihn und geh von der kleineren Zahl so weit nach rechts.\par\smallskip
\grau{a)\ $-0{,}8$ und $-0{,}5$:\quad Abstand $0{,}3$, Hälfte $0{,}15$\quad$\Rightarrow$\quad $-0{,}8 + 0{,}15 = -0{,}65$}\par\medskip
\noindent\begin{minipage}[t]{0.48\linewidth}\raggedright
b)\ $-0{,}4$ und $-0{,}2$
\karo{3}
\end{minipage}\hfill
\begin{minipage}[t]{0.48\linewidth}\raggedright
c)\ $-1$ und $-0{,}6$
\karo{3}
\end{minipage}\par\smallskip
\noindent\begin{minipage}[t]{0.48\linewidth}\raggedright
d)\ $-0{,}3$ und $0{,}1$
\karo{3}
\end{minipage}\par\smallskip
\par{\small Probe: Die Mitte liegt zwischen den beiden Zahlen und ist von beiden gleich weit weg.}
\fusshilfe{\mbox{6:~b) $-0{,}3$; c) $-0{,}8$; d) $-0{,}1$}}
\end{pfaufg}

% L5g-B: Negative Zahlen vergleichen
\begin{pfaufg}{}{7.}{}
Temperaturen am Morgen: Oslo $-3{,}5$\,°C, \ Riga $-3{,}05$\,°C, \ Wien $0{,}5$\,°C, \ Prag $-0{,}8$\,°C, \ Kiew $-4$\,°C. Ordne die Städte von kalt nach warm. Wo war es am kältesten?\karo{3}
\fusshilfe{\mbox{7:~Kiew, Oslo, Riga, Prag, Wien}}
\end{pfaufg}

% L5g-C: Bruch und Prozent: erst umwandeln
\begin{pfaufg}{}{8.}{}
Schreib den Bruch als Dezimalzahl: Erweitere den Nenner auf 10, 100 oder 1000.\par\smallskip
\grau{a)\ $\dfrac{3}{5}$:\quad $\frac{3}{5} = \frac{6}{10} = 0{,}6$}\par\medskip
\noindent\begin{minipage}[t]{0.48\linewidth}\raggedright
b)\ $\dfrac{3}{4}$
\karo{2}
\end{minipage}\hfill
\begin{minipage}[t]{0.48\linewidth}\raggedright
c)\ $\dfrac{7}{20}$
\karo{2}
\end{minipage}\par\smallskip
\noindent\begin{minipage}[t]{0.48\linewidth}\raggedright
d)\ $\dfrac{1}{8}$
\karo{2}
\end{minipage}\hfill
\begin{minipage}[t]{0.48\linewidth}\raggedright
e)\ $\dfrac{9}{25}$
\karo{2}
\end{minipage}\par\smallskip
\fusshilfe{\mbox{8:~b) $0{,}75$; c) $0{,}35$; d) $0{,}125$; e) $0{,}36$}}
\end{pfaufg}

% L5g-C: Bruch und Prozent: erst umwandeln
\begin{pfaufg}{}{9.}{}
Schreib als Dezimalzahl. Prozent heißt Hundertstel: Komma zwei Stellen nach links.\par\smallskip
\grau{a)\ $7\,\%$:\quad $7\,\% = \frac{7}{100} = 0{,}07$}\par\medskip
\noindent\begin{minipage}[t]{0.48\linewidth}\raggedright
b)\ $45\,\%$
\karo{2}
\end{minipage}\hfill
\begin{minipage}[t]{0.48\linewidth}\raggedright
c)\ $4{,}5\,\%$
\karo{2}
\end{minipage}\par\smallskip
\noindent\begin{minipage}[t]{0.48\linewidth}\raggedright
d)\ $120\,\%$
\karo{2}
\end{minipage}\hfill
\begin{minipage}[t]{0.48\linewidth}\raggedright
e)\ $0{,}5\,\%$
\karo{2}
\end{minipage}\par\smallskip
\fusshilfe{\mbox{9:~b) $0{,}45$; c) $0{,}045$; d) $1{,}2$; e) $0{,}005$}}
\end{pfaufg}

% L5g-C: Bruch und Prozent: erst umwandeln
\begin{pfaufg}{}{10.}{}
Wandle erst in Dezimalzahlen um, dann vergleiche. Setze $<$, $>$ oder $=$.\par\smallskip
\grau{a)\ $\dfrac{3}{4}$ \fbox{\phantom{$<$}} $0{,}8$:\quad $\frac{3}{4} = 0{,}75$\quad$\Rightarrow$\quad $0{,}75 < 0{,}80$, also $\frac{3}{4} < 0{,}8$}\par\medskip
\noindent\begin{minipage}[t]{0.48\linewidth}\raggedright
b)\ $6\,\%$ \fbox{\phantom{$<$}} $0{,}6$
\karo{3}
\end{minipage}\hfill
\begin{minipage}[t]{0.48\linewidth}\raggedright
c)\ $0{,}35$ \fbox{\phantom{$<$}} $\dfrac{7}{20}$
\karo{3}
\end{minipage}\par\smallskip
\noindent\begin{minipage}[t]{0.48\linewidth}\raggedright
d)\ $\dfrac{2}{5}$ \fbox{\phantom{$<$}} $45\,\%$
\karo{3}
\end{minipage}\hfill
\begin{minipage}[t]{0.48\linewidth}\raggedright
e)\ $\dfrac{1}{8}$ \fbox{\phantom{$<$}} $0{,}13$
\karo{3}
\end{minipage}\par\smallskip
\fusshilfe{\mbox{10:~b) $6\,\% < 0{,}6$; c) $0{,}35 = \frac{7}{20}$; d) $\frac{2}{5} < 45\,\%$; e) $\frac{1}{8} < 0{,}13$}}
\end{pfaufg}

% L5g-D: Quadrat und Wurzel
\begin{pfaufg}{}{11.}{}
Rechne das Quadrat aus: die Zahl mal sich selbst – nicht mal 2.\par\smallskip
\grau{a)\ $0{,}4^2$:\quad $0{,}4^2 = 0{,}4 \cdot 0{,}4 = 0{,}16$}\par\medskip
\noindent\begin{minipage}[t]{0.48\linewidth}\raggedright
b)\ $0{,}2^2$
\karo{2}
\end{minipage}\hfill
\begin{minipage}[t]{0.48\linewidth}\raggedright
c)\ $0{,}5^2$
\karo{2}
\end{minipage}\par\smallskip
\noindent\begin{minipage}[t]{0.48\linewidth}\raggedright
d)\ $1{,}1^2$
\karo{2}
\end{minipage}\hfill
\begin{minipage}[t]{0.48\linewidth}\raggedright
e)\ $0{,}1^2$
\karo{2}
\end{minipage}\par\smallskip
\fusshilfe{\mbox{11:~b) $0{,}04$; c) $0{,}25$; d) $1{,}21$; e) $0{,}01$}}
\end{pfaufg}

% L5g-D: Quadrat und Wurzel
\begin{pfaufg}{}{12.}{}
Zieh die Wurzel: Welche Zahl mal sich selbst gibt die Zahl unter der Wurzel? Prüfe mit dem Quadrat.\par\smallskip
\grau{a)\ $\sqrt{0{,}25}$:\quad $0{,}5 \cdot 0{,}5 = 0{,}25$, also $\sqrt{0{,}25} = 0{,}5$}\par\medskip
\noindent\begin{minipage}[t]{0.48\linewidth}\raggedright
b)\ $\sqrt{0{,}09}$
\karo{2}
\end{minipage}\hfill
\begin{minipage}[t]{0.48\linewidth}\raggedright
c)\ $\sqrt{1{,}44}$
\karo{2}
\end{minipage}\par\smallskip
\noindent\begin{minipage}[t]{0.48\linewidth}\raggedright
d)\ $\sqrt{0{,}04}$
\karo{2}
\end{minipage}\hfill
\begin{minipage}[t]{0.48\linewidth}\raggedright
e)\ $\sqrt{0{,}49}$
\karo{2}
\end{minipage}\par\smallskip
\fusshilfe{\mbox{12:~b) $0{,}3$; c) $1{,}2$; d) $0{,}2$; e) $0{,}7$}}
\end{pfaufg}

% L5g-D: Quadrat und Wurzel
\begin{pfaufg}{}{13.}{}
Die Wurzel geht nicht auf. Zwischen welchen zwei Zehnteln liegt sie? Probiere Quadrate, bis die Zahl dazwischen liegt.\par\smallskip
\grau{a)\ $\sqrt{2}$:\quad $1{,}4^2 = 1{,}96$, \ $1{,}5^2 = 2{,}25$\quad$\Rightarrow$\quad $1{,}4 < \sqrt{2} < 1{,}5$}\par\medskip
\noindent\begin{minipage}[t]{0.48\linewidth}\raggedright
b)\ $\sqrt{3}$
\karo{3}
\end{minipage}\hfill
\begin{minipage}[t]{0.48\linewidth}\raggedright
c)\ $\sqrt{5}$
\karo{3}
\end{minipage}\par\smallskip
\noindent\begin{minipage}[t]{0.48\linewidth}\raggedright
d)\ $\sqrt{0{,}5}$
\karo{3}
\end{minipage}\par\smallskip
\fusshilfe{\mbox{13:~b) $1{,}7 < \sqrt{3} < 1{,}8$; c) $2{,}2 < \sqrt{5} < 2{,}3$; d) $0{,}7 < \sqrt{0{,}5} < 0{,}8$}}
\end{pfaufg}

% L5g-E: Alles gemischt: vergleichen und entscheiden
\begin{pfaufg}{}{14.}{}
\grau{a)\ Welcher Wert ist der kleinste? Kreuze an. \kreuz{$0{,}7$} \kreuz{$0{,}1^2$} \kreuz{$3\,\%$} \kreuz{$0{,}025$}\par\smallskip $0{,}7 = 0{,}700$; \ $0{,}1^2 = 0{,}010$; \ $3\,\% = 0{,}030$; \ $0{,}025$\par Zehntel: nur $0{,}700$ hat welche. Hundertstel: $0{,}010$ hat 1, die anderen 2 und 3.\par Der kleinste ist $0{,}1^2 = 0{,}01$.}
\par\medskip
b)\ Welcher Wert ist der kleinste? Kreuze an.\par\smallskip \gkreuz{$0{,}9$}\gkreuz{$0{,}3^2$}\par \gkreuz{$8\,\%$}\gkreuz{$0{,}085$}
\fusshilfe{\mbox{14:~b) $8\,\% = 0{,}08$}}
\end{pfaufg}

% L5g-E: Alles gemischt: vergleichen und entscheiden
\begin{pfaufg}{}{15.}{}
Ordne von klein nach groß: \ $0{,}3$; \ $-\frac{3}{4}$; \ $\sqrt{0{,}16}$; \ $-0{,}7$.\karo{3}
\fusshilfe{\mbox{15:~$-\frac{3}{4} < -0{,}7 < 0{,}3 < \sqrt{0{,}16}$}}
\end{pfaufg}

% L5g-E: Alles gemischt: vergleichen und entscheiden
\begin{pfaufg}{}{16.}{}
Nur eine Aussage ist wahr. Kreuze sie an.\par\medskip\noindent\hspace*{7mm}\mbox{$\square$\ $\sqrt{2} > \dfrac{3}{2}$}\hspace{10mm}\mbox{$\square$\ $0{,}3^2 > 0{,}3$}\hspace{10mm}\mbox{$\square$\ $-0{,}5 < -\dfrac{3}{5}$}\hspace{10mm}\mbox{$\square$\ $\dfrac{5}{8} > 0{,}6$}\rule[-3mm]{0pt}{8mm}\par\smallskip 
\fusshilfe{\mbox{16:~$\frac{5}{8} > 0{,}6$}}
\end{pfaufg}

% L5g-E: Alles gemischt: vergleichen und entscheiden
\begin{pfaufg}{}{17.}{}
Drei Läden haben dieselbe Jacke. Laden A: „Ein Viertel billiger!“ \ Laden B: „22\,\% Rabatt“ \ Laden C: „Jeden fünften Euro schenken wir dir!“ \ Wo ist der Rabatt am größten? Um wie viele Prozentpunkte ist er größer als beim zweitbesten Laden?\karo{3}
\fusshilfe{\mbox{17:~Laden A; 3 Prozentpunkte mehr als B}}
\end{pfaufg}

% L5g-E: Alles gemischt: vergleichen und entscheiden
\begin{pfaufg}{}{18.}{}
Mia taucht in $4{,}8$\,m Tiefe, ihr Bruder in $2{,}4$\,m Tiefe. Genau in der Mitte zwischen beiden schwimmt ein Fisch. Ein Foto gelingt nur tiefer als $3{,}5$\,m. Kann Mia den Fisch fotografieren?\karo{3}
\fusshilfe{\mbox{18:~Ja; Fisch in $3{,}6$\,m Tiefe}}
\end{pfaufg}

\par\vfill\noindent{\footnotesize Vorher: Vergleichen, Ordnen, Runden (Kl. 6, HMF)\quad\textbullet\quad Weiter: Potenzen und Wurzeln\par}
\end{document}
